\subsubsection*{引理 4.5.8 (Dwork 迹引理).}
设 $b\in\mathbb Q$, $b>0$, $k\in\mathbb Z$, $k\geq2$, 且 $G\in L(b/k)$.
则复合映射

\[
L(b)\xrightarrow{\ \text{限制}\ }L(b/k)
\xrightarrow{\ \times G\ }L(b/k)
\xrightarrow{\ \psi_k\ }L(b),
\]

记为 $\psi_k\circ G$, 是全连续的, 其迹 (参见 [4]) 由下式给出

\begin{equation}\tag{4.5.9}
\operatorname{tr}(\psi_k\circ G)
=\frac1{(k-1)^{n+3}}
\sum_{X\in(\mu_{k-1}(\Omega))^{n+3}}G(X).
\end{equation}

\begin{editorialnote}
原书第421页及英文底本在 (4.5.9) 的分母写作 $(k-1)^{n+2}$, 但求和遍及 $n+3$ 个坐标.
取 $G=1$ 时, $\psi_k$ 的迹为 $1$, 根之和有 $(k-1)^{n+3}$ 项, 因而分母应为 $(k-1)^{n+3}$.
这也与下一式 (4.5.10) 的因子相符.
\end{editorialnote}

将此引理用于

\[
\left\{
\begin{aligned}
k&=q^\nu,\\
G(X)&=\prod_{j=0}^{\operatorname{ord}(q^\nu)-1}
H^{\tau^j}(X^{p^j})
\in L\!\left(\frac{p}{(p-1)q^\nu}\right),\\
b&=\frac1{p-1}.
\end{aligned}
\right.
\]

公式 (4.4.6) 化为

\begin{equation}\tag{4.5.10}
\begin{aligned}
(q^\nu-1)\operatorname{card}V^*(\mathbb F_{q^\nu})
&=\frac{(q^\nu-1)^{n+2}}{q^\nu}\\
&\quad+\frac{(q^\nu-1)^{n+3}}{q^\nu}
\operatorname{tr}\!\left(
\left.\psi_{q^\nu}\circ
\prod_{j=0}^{\operatorname{ord}(q^\nu)-1}H^{\tau^j}(X^{p^j})
\right|L\!\left(\frac1{p-1}\right)\right).
\end{aligned}
\end{equation}

设 $L^0(b)$ 是 $L(b)$ 中“在零点消失”的元素组成的理想, 即不含常数项的元素组成的理想.
它在 $\psi_k\circ G$ 下稳定, 且显然

\begin{equation}\tag{4.5.11}
\operatorname{tr}(\psi_k\circ G)
=G(0)+\operatorname{tr}\!\left(\left.\psi_k\circ G\right|L^0(b)\right).
\end{equation}

由于 $H(0)=1$, 公式 (4.5.10) 可改写为

\begin{equation}\tag{4.5.12}
\begin{aligned}
\operatorname{card}V^*(\mathbb F_{q^\nu})
&=(q^\nu-1)^{n+1}\\
&\quad+\frac{(q^\nu-1)^{n+2}}{q^\nu}
\operatorname{tr}\!\left(
\left.\psi_{q^\nu}\circ
\prod_{j=0}^{\operatorname{ord}(q^\nu)-1}H^{\tau^j}(X^{p^j})
\right|L^0\!\left(\frac1{p-1}\right)\right).
\end{aligned}
\end{equation}
